\section{Mobile App Testing}\label{sec:MobileAppTesting}

\subsection{Unit Testing}
Unit testing is the process of verifying the correctness of a single function,
method, or class in isolation. It is considered the foundation of the testing
pyramid. The primary goal is to ensure that the business logic and algorithms
function as expected under a variety of conditions without any interaction with
the User Interface (UI) or external dependencies such as databases or web
servers \cite{flutterunittest}.

As described by Torres \cite{torres2024testing}, unit tests help ensure that
specific pieces of logic produce the expected output given a particular input.
These tests are characterized by their execution speed and low maintenance
cost. In Flutter, unit tests typically utilize the \texttt{test} package to
define expectations and assertions for Dart classes.

\subsection{Widget Testing}
Widget testing, often referred to as component testing in other frameworks,
sits between unit and integration testing. A widget test verifies the behavior
of a single widget or a small group of widgets \cite{flutterwidgettest}.

According to Bharti \cite{bharti2024widget}, the objective of widget testing is
to ensure UI reliability and quality by asserting that the widget's visual
layout is rendered correctly and that it responds appropriately to user
interactions, such as button taps or text entry.

Unlike unit tests, widget tests require a simulated environment. The Flutter
framework provides the \texttt{WidgetTester} utility, which allows developers
to build and interact with widgets in a test environment that is significantly
faster than a full UI system but complex enough to render the widget tree
\cite{flutterwidgettest}.

\subsection{Integration Testing}
Integration testing involves testing the application as a complete system. This
methodology validates that all individual components—including the UI, business
logic, and external services—work together correctly \cite{torres2024testing}.

These tests simulate a real user's journey through the application, such as
logging in, navigating through screens, and performing data updates. Because
integration tests typically run on real devices or emulators to capture
accurate performance data and visual rendering, they are generally slower and
more resource-intensive than unit or widget tests.
\subsection{Testing Architecture and Tools}
We utilized the following standard Flutter testing libraries:
\begin{itemize}
    \item \textbf{flutter\_test}: The core framework used to write and run assertions.\cite{flutter_testing_docs}
    \item \textbf{mockito}: Used to simulate external dependencies. Since our ViewModels depend on Repositories (e.g., \texttt{BookingRepository}) and Plugins (e.g., \texttt{ImagePicker}), we created "Mock" objects to return controlled dummy data during tests.\cite{mockito_pub}
    \item \textbf{build\_runner}: Used to automatically generate the Mock classes \newline(e.g., \texttt{MockMajorEventRepository}) using the \texttt{@GenerateMocks} annotation, ensuring our test mocks stay up-to-date with our actual interfaces.\cite{build_runner_pub}
    \item \textbf{WidgetTester}: Provided by the \texttt{flutter\_test} package, this allows us to build and interact with widgets in a test environment without running the full app on an emulator.\cite{flutter_widgettester}
    \item \textbf{network\_image\_mock}: Since widget tests run in a restrictive environment without real internet access, standard \texttt{Image.network} widgets would fail. We used this package to wrap our tests and mock successful image loads.\cite{tigerasks2024provider, network_image_mock}
    \item \textbf{Provider Injection}: Many of our widgets (e.g., \texttt{DoctorCard}) rely on \newline\texttt{ChangeNotifierProvider}. In tests, we injected "Mock Providers" (created with Mockito) to supply controlled data and verify that buttons trigger the correct business logic.\cite{tigerasks2024provider}
\end{itemize}

\subsection{Unit Testing Implementation}
To ensure the reliability of the application's business logic and maintain a
stable codebase, we implemented a comprehensive Unit Testing strategy. This
strategy focused on verifying the behavior of our \textbf{ViewModels} in isolation from the UI and external dependencies (such as the
backend API or local device storage).

\subsubsection{Booking Logic (BookingProvider)}
The booking process involves complex state management where selecting a doctor
resets previously selected dates to prevent invalid slots.
\begin{itemize}
    \item \textbf{Scenario Tested:} We verified that selecting a new doctor automatically clears the \texttt{selectedDate} and \texttt{selectedTime} to enforce valid scheduling.
    \item \textbf{Optimistic Updates:} We tested the \texttt{confirmBooking} method to ensure that once the repository returns success, the local appointment list updates immediately without requiring a full page refresh.
\end{itemize}

\begin{verbatim}
// Code Snippet: Verifying Booking Reset Logic (booking_provider_test.dart)
test('Selecting a doctor should reset Date and Time', () {
  // Arrange: Set a date first
  provider.setTimeSlot(DateTime(2026, 1, 1), '10:00 AM');
  
  // Act: Select a new doctor
  provider.selectDoctor(createDummyDoctor());

  // Assert: Doctor is set, but time/date are wiped (null)
  expect(provider.selectedDoctor, isNotNull);
  expect(provider.selectedDate, null);
  expect(provider.selectedTime, null);
});
\end{verbatim}

\subsubsection{Document Uploads (DocumentAddViewModel)}
This module interacts with native platform channels (File Picker, Image
Picker). Testing it requires intercepting these platform calls.
\begin{itemize}
    \item \textbf{Validation Logic:} We verified that the \texttt{isValid} getter returns \texttt{false} unless both a "Title" and a "File" are present, preventing users from submitting incomplete forms.
    \item \textbf{Platform Mocking:} Since Unit Tests run on a computer (not a phone), calls to \texttt{getApplicationDocumentsDirectory} would normally crash. We used \texttt{MethodChannel} mocking to intercept these calls and return a temporary directory on the test machine.
\end{itemize}

\subsubsection{Medical History Timeline (MajorEventProvider)}
This provider handles fetching the patient's "Major Events" and their
sub-nodes.
\begin{itemize}
    \item \textbf{State Management:} We tested the \texttt{fetchEvents} method to ensure the \texttt{isLoading} boolean correctly toggles to \texttt{true} when the call starts and \texttt{false} when data is received.
    \item \textbf{Data Clearing:} We verified that fetching details for a new event (e.g., "Surgery") correctly clears the old data from the previous event to prevent showing stale information.
\end{itemize}

\begin{verbatim}
// Code Snippet: Mocking Repository Responses (major_event_provider_test.dart)
test('fetchEvents success updates list and loading state', () async {
  // Arrange: Stub the repo to return a dummy list immediately
  when(mockRepo.getMajorEvents())
      .thenAnswer((_) async => [createDummyEvent()]);

  // Act: Trigger the fetch
  final future = provider.fetchEvents();
  
  // Assert: Loading should be true while fetching
  expect(provider.isLoadingEvents, true);

  await future;

  // Assert: Data is populated and loading is false
  expect(provider.events.length, 1);
  expect(provider.isLoadingEvents, false);
});
\end{verbatim}

\subsubsection{Test Coverage Results}
By isolating the ViewModels and mocking the repositories, we achieved high test
coverage for the application's core logic. This approach allowed us to:
\begin{enumerate}
    \item \textbf{Identify Edge Cases:} We caught issues where invalid form submissions were possible before implementing the validation checks.
    \item \textbf{Refactor Safely:} We could optimize the \texttt{BookingProvider} logic with confidence, knowing the tests would fail if we broke the "reset date" functionality.
    \item \textbf{Improve Efficiency:} Running these unit tests takes less than 2 seconds, compared to manual testing which would require navigating through the entire app.
\end{enumerate}
